
\documentclass[letter, 11pt]{article}

    %%%%%%%%%%%%%%%%%%%%%%%%%%%%%%
    \oddsidemargin  0.0in
    \evensidemargin 0.0in
    \textwidth      6.8in
    \headheight     0.0in
    \topmargin      0.0in
    \textheight=8.4in
    \parindent=0in
    %%%%%%%%%%%%%%%%%%%%%%%%%%%%%%

\include{MathCommands}

\usepackage{amsthm,amsfonts,amssymb,amsmath}
\usepackage{graphicx}   %To include graphics
\usepackage[small,compact]{titlesec}

\usepackage[utf8]{inputenc}

\usepackage[francais]{babel}

\usepackage{enumerate}
\usepackage{fancyhdr}
\pagestyle{fancy}

\renewcommand{\headrulewidth}{1pt}
\fancyhead[C]{Automne 2021}
\fancyhead[L]{STT-2100 Régression}
\fancyhead[R]{Audrey-Anne Vallée}

\newcommand{\var}{\mathbb{V}\text{ar}}
\newcommand{\cov}{\mathbb{C}\text{ov}}
\newcommand{\esp}{\mathbb{E}}

\begin{document}

\vspace*{0.05 in}
\begin{center}
{\Large Devoir 1 \\
Dû le 1er octobre 2021 à 17h00}
\end{center}
\vspace*{0.1 in}


À faire individuellement ou en équipe de deux. \\

Vous devez remettre \underline{trois fichiers} par équipe:
\begin{itemize}
\item Un fichier PDF nommé {\tt EquipeX-1234.pdf} avec vos réponses aux questions 1 à 4. Vous pouvez simplement numériser une copie papier écrite à la main. Remplacez {\tt X} par le numéro de votre équipe.
\item Un fichier PDF nommé {\tt EquipeX-5.pdf} avec vos réponses de la question 5. Celui-ci doit être rédigé à l'ordinateur et être d'une longueur maximale de 3 pages. %N'indiquez que le numéro de votre équipe sur ce document; certains rapports seront partagés avec les autres étudiants du cours après la remise du devoir.
\item Un document {\tt EquipeX.R}, ou {\tt EquipeX.sas} avec les codes utilisés pour la question 5.\\
\end{itemize}


\underline{La qualité de la langue sera évaluée}, incluant l'orthographe, la grammaire et la syntaxe. Un maximum de 10 points pourra être retranché pour mauvaise qualité du français. \\

Un travail remis en retard ne sera pas accepté. Les étudiants remettant un devoir en retard recevront la note 0.
%Un travail remis entre 17h00 le 15 octobre et 17h00 le 16 octobre sera pénalisé de 20 points. %Un travail remis entre 17h00 le 8 octobre et 17h00 le 9 octobre sera pénalisé de 30 points.
%Aucun travail ne sera accepté après cette date.

\vspace{0.5cm}

\underline{\hspace{17cm}}\\

\textbf{Question 1 (20 pts)} \\

Soit le modèle de régression
$$
Y_i = \beta_0\,+\,\beta_1\, x_i \,+\, \varepsilon_i\qquad \text{ pour } i=1,\ldots,n
$$
tel que 
$$ \sum_{i = 1}^n x_i = 0.$$
On fait l'hypothèse que les erreurs $\varepsilon_i$ sont indépendantes et distribuées selon la loi normale avec espérance~0 et variance $\sigma^2$. 

\begin{enumerate}[(a)]
%\item[(a)] (3 pts) Donnez la distribution de $Y_i$. Justifiez votre réponse.
\item  (9 pts) Calculez les estimateurs des moindres carrés de $\beta_0$ et de $\beta_1$. {\bf Vous devez exprimer les estimateurs sous leur forme la plus réduite.} Justifiez votre démarche. (Vous pouvez utiliser directement tout résultat montré dans cours.)
\item (8 pts) {\bf Sans utiliser les résultats du cours}, calculez l'espérance et la variance de $\hat\beta_0$. Justifiez votre démarche et vos calculs.
\item (8 pts)  {\bf Sans utiliser les résultats du cours}, calculez l'espérance et la variance de $\hat\beta_1$. Justifiez votre démarche et vos calculs.
\end{enumerate}

%\vspace{0.2 in}
\newpage
\textbf{Question 2 (15 pts)} \\

Soit le modèle de régression
$$
Y_i = \beta_0\,+\,\beta_1\, x_i \,+\, \varepsilon_i\qquad \text{ pour } i=1,\ldots,n
$$
tel que  $\mathrm{E}\left(\varepsilon_i\right) = 0$, $\mathrm{E}\left(\varepsilon_i^2\right) = \sigma^2$ et $\mathrm{E}\left(\varepsilon_i \varepsilon_j \right) = 0$ pour $i\neq j$. On estime $\beta_0$ et $\beta_1$ à l'aide de la méthode des moindes carrés et on obtient $\widehat{\beta}_0$ et $\widehat{\beta}_1$. Les valeurs de la variable $x$ sont en pouces et on aimerait écrire le modèle avec la variable explicative $z_i$ en centimètres. On suppose qu'un pouce est égal à $c$ centimètres (avec un $c$ connu et fixe). On veut étudier le modèle
$$
Y_i = \beta_0^* \,+\,\beta_1^* \, z_i \,+\, \varepsilon_i\qquad \text{ pour } i=1,\ldots,n.
$$

\vspace{0.2 in}
\begin{itemize}
\item[(a)] (10 pts) Calculez les estimateurs de $\beta_0^*$ et $\beta_1^*$ et exprimez les en fonction de $\widehat{\beta}_0$ et $\widehat{\beta}_1$. Justifiez votre réponse.
\vspace{0.1 in}
\item[(b)] (5 pts) Montrez que la valeur de $R^2$ est la même pour les deux modèles.
\end{itemize}

\vspace{0.5cm}

\underline{\hspace{17cm}}\\


\textbf{Question 3 (10 pts)} \\

Soit le modèle 
$$
Y_i = \mu \,+\, \varepsilon_i\qquad \text{ pour } i=1,\ldots,n
$$
tel que  $\mathrm{E}\left(\varepsilon_i\right) = 0$, $\mathrm{E}\left(\varepsilon_i^2\right) = \sigma^2$ et $\mathrm{E}\left(\varepsilon_i \varepsilon_j \right) = 0$ pour $i\neq j$.
Sans utiliser de résultats du cours, calculez l'estimateur des moindres carrés de~$\mu$. Justifiez votre réponse.

\vspace{0.5cm}

\underline{\hspace{17cm}}\\

\textbf{Question 4 (10 pts)} \\

Considérons le modèle de régression linéaire simple usuel avec erreurs normales tel que vu en cours. On a déjà déterminé la distribution de $\hat{\beta}_0$ et celle de $\hat{\beta}_1$. On s'intéresse à l'estimation de $\sigma^2$. \\

Montrer que 
$$\hat\sigma_{MV}^2=\frac{1}{n}\sum\limits_{i=1}^n \hat\varepsilon_i^2$$ 
est un estimateur  biaisé de la variance autour de la droite $\sigma^2$. (Vous pouvez utiliser directement  la distribution des résidus abordée dans le chapitre 4.)




\newpage

\vspace{0.1 in}
\textbf{Question 5 (40 pts)} \\

On s'intéresse à la relation entre le poids et la taille des poissons d'un lac. On dispose d'un échantillon de poissons pour lesquels on a mesuré le poids (en grammes) et la taille (en centimètres).
On vous demande de trouver un modèle linéaire pour prédire le poids d'un poisson à partir de sa taille.  
Chaque équipe recevra un jeu de données avec les mesures prises pour un échantillon de poissons différent. %, sélectionné au hasard parmi un plus grand jeu de données.
Pour chaque poisson, on vous donne la taille (en cm) et le poids (en gr) du poisson.  
Vous devez rédiger des parties d'un rapport d'analyse de régression en répondant aux questions suivantes. Utilisez $\alpha = 0,01$ pour toutes vos inférences.\\ 

{\bf Dans le fichier PDF}: identifiez clairement la question associée à votre réponse (par exemple (i) a)). Donnez les réponses aux questions et les tableaux de résultats obtenus avec votre logiciel qui sont nécessaires pour
répondre aux questions. Les figures et tableaux inclus dans vos réponses doivent être présentés de façon professionnelle. Ils doivent avoir des titres clairs et pertinents. Chaque figure présentée doit être discutée dans vos réponses. \\

\textbf{Sorties statistiques:}
Vous devez fournir les programmes R ou SAS utilisés pour obtenir vos réponses. Un programme qui ne fonctionne pas ou qui est mal annoté entraînera une pénalité. \\


\begin{enumerate}[(i)]
	\item {\bf Décription du jeu de données et de l'objectif poursuivi.}
\begin{enumerate}[a)]
\item Identifiez et définissez les variables qui nous intéressent.
%\item Mentionner le nombre et la source des observations.
\item Décrivez {\it brièvement} la distribution de chacune des variables (graphiques et statistiques descriptives). Commentez la distribution conjointe des variables à l'aide d'un diagramme de dispersion.  (8 lignes maximum, sans compter les figures).
\item Décrire l'objectif de l'analyse de régression. (4 lignes maximum).
\end{enumerate}
	\item {\bf Choix de modèle utilisé.}
\begin{enumerate}[a)]
	%\item Si on choisit de transformer les variables avant de procéder \`a l'analyse, justifier son choix (raisons théoriques, expérience antérieure avec ce type de données, intuition seulement\,?).
	%\item Au besoin, simplifier le mod\`ele (par exemple, en refaisant la régression sans $\beta_0$).
	\item Vérifiez l'adéquation du mod\`ele RLS et le respect des postulats au moyen de graphiques et tests pertinents. Modifiez le modèle de façon à ce que vous soyez satisfait de la justesse des postulats. (On ne peut pas toujours modifier les variables pour que tous les postulats soient exactement vérifiés;  choisissez alors le meilleur modèle possible). Dans le fichier réponse, formulez clairement le modèle et les postulats du mod\`ele retenu (i.e. écrire le modèle complet en language mathématique). Présentez seulement le modèle final, et non tous les modèles ajustés. Motivez le choix de modèle final (8 lignes maximum). Incluez toutes les étapes de votre démarche (différents modèles, vérification des postulats) dans votre programme R ou SAS. 
\end{enumerate}
	\item {\bf Présentation des résultats.}
\begin{enumerate}
	\item Donnez les estimations des param\`etres (${\hat \beta}_0$, ${\hat \beta}_1$, ${\hat \sigma}^2$).
	\item Fournissez des intervalles de confiance pour les paramètres (${\hat \beta}_0$, ${\hat \beta}_1$) et une interprétation en contexte. (8 lignes maximum pour l'interprétation).
	\item Discutez de la signification statistique des coefficients ${\hat \beta}_0$ et ${\hat \beta}_1$ à l'aide de tests $t$. (6 lignes maximum).
	\item L'analyse de variance permet de faire un test d'hypothèses. Donnez les hypothèses de ce test et discutez de sa signification statistique pour vos données. (4 lignes maximum pour la discussion).
\end{enumerate}
\end{enumerate}




\end{document}




\vspace{0.3in}
\textbf{Question 4 (10 pts)} (Design d'expérience)\\

On s'intéresse à un médicament sensé réduire le nombre de crises chez les individus atteints d'epilepsie. Vous devez aider à planifier une expérience pour étudier l'effet de la dose du médicament sur le nombre de crises. Puisque vous vous intéressez à la relation entre deux variables quantitatives, vous planifiez analyser les données avec un modèle de régression linéaire simple. \\

Soit $x_1, \ldots, x_n$ les valeurs de la dose pour les 20 individus de votre étude.\\ Pour chacun, vous pouvez attribuer n'importe quelle dose entre 2 mg et 21 mg. \\

\begin{itemize}
%\item [(a)] Pour le modèle de régression linéaire simple $Y_i = \beta_0 + \beta_1 x_i + \epsilon_i$, prouvez que la variance de $\beta_0$ est minimale si les $x_i$ sont tels que $\bar{x}$ = 0.
\item [(a)] (5 pts) Quelle est la meilleure fa\c{c}on de choisir les valeurs $x_1,\ldots,x_n$ pour minimiser la variance de votre estimateur $\hat{\beta_1}$? Justifiez votre réponse.
\vspace{0.1 in}
\item [(b)] (5 pts) Pourquoi ne serait-ce pas vraiment une bonne idée de choisir les valeurs de $x_i$ de cette façon, même si cela minimise la variance de votre estimateur?
\end{itemize}


\end{document} 